\begin{textsample}{NEG Dim 2 – vem\_ed\_13 – Score 4.00 – vem\_ed\_13/t056.txt}  \label{ex:f2_neg_111}
NOVAS COLABORAÇÕES COM PORTUGAL A equipe dos PCIs / Cesu começou, em 2022, uma colaboração com o Instituto Superior D.
Dinis (ISDOM), em Portugal.
É a terceira Instituição de Ensino Superior portuguesa a elaborar PCIs com Fatecs: desde 2021, a Universidade de Aveiro e o Instituto Politécnico de Viseu realizam projetos com Fatecs, conforme noticiado, respectivamente, nas edições 12 e 6 de VEm: https://cesu.cps.sp.gov.br/vem-newsletter-bimestral-dos-projetos-colaborativos-internacionais-pci/ Com os docentes do ISDOM e seus estudantes, estão trabalhando as Fatecs das cidades de Cruzeiro, Guaratinguetá, Lins, São Caetano do Sul e Sertãozinho.
As unidades localizadas em Cruzeiro, Guaratinguetá e Sertãozinho fazem estudos em gestão da produção industrial.
Eis alguns dos temas abordados: planejamento e controle de serigrafia em garrafas de vidro; mapeamento dos processos de fabricação de um molde (e de seus sistemas de gerenciamento) e alteração de layout para melhoria no fluxo de produção da indústria moveleira.
Atuam com suas turmas Aníbal Evaristo Fernandes (coordenador do curso de Gestão da Produção Industrial / GPI), Leônidas Morais (Fatec Cruzeiro), Juliano Endrigo Sordan (coordenador do curso de GPI da Fatec Sertãozinho) e Luiz Antônio Alvarenga (coordenador do curso de GPI da Fatec Guaratinguetá), com a professora Arminda Pata (ISDOM).
Na Fatec de Guaratinguetá, existe mais um projeto, sobre processos, manufatura, embalagem e desenvolvimento de produtos em polímeros e vidros.
Os professores Rosinei Batista Ribeiro, Ana Raquel Elisa Satim Rodrigues (Fatec Guaratinguetá), Susana Ramalho, Filipa Cardal e Antônio Guilherme Cristiano (ISDOM) orientam quatro grupos de estudantes das duas instituições. “Estudo comparativo sobre gestão dos resíduos sólidos em empresas brasileiras e portuguesas” é o tema do PCI conduzido por Sandro da Silva Pinto (coordenador do curso de Gestão da Qualidade da Fatec Lins) e Fernando Frachone Neves (coordenador do curso de Gestão Empresarial da Fatec Sertãozinho), com Arminda Pata e José Mendes (ISDOM).
Os docentes distribuíram informações sobre políticas públicas e legislação relacionadas ao tema nos dois países, para que os estudantes verifiquem o cumprimento dessas leis e diretrizes.
Paula Daniela da Silva Monciatti, professora do curso de Comércio Exterior da Fatec São Caetano do Sul, pesquisa com Fernando Remondes (ISDOM) e seus alunos o fluxo de materiais e a gestão de estoques da cadeia do vidro no setor metal-mecânico.

% matched lemmas: 
\end{textsample}
